\section{Conclusions}
\label{sec:conclusions}

Cap dels encoders que hem provat (CNN, GRU, BiLSTM més profundes, Transformer o híbrids) supera la BiLSTM d'una
capa del baseline més enllà de la variabilitat entre llavors. Amb frases curtes i 4.000 exemples d'entrenament,
aquesta xarxa ja té més capacitat de la que les dades poden aprofitar. Les millores han vingut, en canvi, de la
representació d'entrada: dels vectors GloVe, de la char-CNN i del dropout a l'embedding. Combinats, porten la F1 de
test de 0,946 a 0,955 i la F1 macro de 0,727 a 0,777, i la millora es manté quan deixem entrenar tots dos models
sense \mindelta.

Una part dels resultats ha contradit les nostres hipòtesis, i en bona part és el que més ens ha ensenyat: el
Transformer millora amb més capes, el dropout no fa res a la BiLSTM gran i el CRF, sense restriccions explícites,
no elimina les seqüències impossibles. També hem après que el \mindelta{} que vam afegir era massa estricte. Als
dos models on ho hem comprovat ens costava entre 0,2 i 0,4 punts, i si tornéssim a fer l'estudi no el faríem
servir. Els pesos de classe, per la seva banda, mostren fins a quin punt la mètrica de selecció condiciona el
resultat final, ja que milloren les classes rares però, amb la F1 micro com a criteri, no entren al model.

L'anàlisi d'errors apunta cap a on continuar. La majoria dels errors que queden són confusions de rol que depenen
de la intenció de la frase, i un model que aprengués alhora la intenció (la primera part de la pràctica) i les
entitats disposaria d'aquesta informació. També provaríem la decodificació restringida al CRF.
